\chapter*{Abstract}
\addcontentsline{toc}{chapter}{Abstract}

%% Edit below this line
This project investigates training-free Video Anomaly Detection (VAD) in production and
assembly lines using Vision-Language Models (VLMs), avoiding any task-specific training.
The concrete scene studied is a part travelling on a conveyor. Three method families are studied: direct prompting, which establishes
baselines and reveals the model's bias toward predicting ``normal''; a multi-persona
majority-voting scheme on a locally deployed open-weight model
(Qwen3-VL-8B-Instruct), in which diverse inspector personas vote on a verdict and an
\emph{articulate-first} output format (describe-then-decide) suppresses systematic
false negatives; and \emph{VERA-lite}, a guiding-question optimisation pipeline with a
blind evaluation protocol, alternative input representations, reasoning-budget control,
and a two-view ensemble. On a purpose-built $13$-clip pen-conveyor dataset, the ensemble
of a continuous-video view and a discrete dense-frame view reaches $13/13$ accuracy with
zero false positives by combining their disjoint blind spots, whereas no single
representation succeeds, exposing an intrinsic continuous-versus-discrete trade-off for
temporal anomalies. Repeated runs show the true blind accuracy of a single view is
$0.846$ ($11/13$, the same in all three runs) rather than a lucky single-run $0.92$, and
that capping the model's reasoning budget yields a $\sim$41\% latency reduction at
identical accuracy while also making the per-clip verdicts fully deterministic. We report
negative results honestly and discuss limitations and future work, including larger
datasets, computer-vision baselines, and frame-level evaluation.

%% Until here
\vfill
%% Edit after {Keywords:}
\textbf{Keywords:} video anomaly detection, vision-language models, zero-shot,
prompting, ensemble, industrial inspection.
\clearpage
